\section{The two-stage observer schedule (Proposition 6.4(b))}\label{S:observer}

Information structure: the blind cell $\{1,2,3\}$ on $[0,\tau_1)$, the partition $\{1\}\,|\,\{2,3\}$ on $[\tau_1,\tau)$, full revelation from $\tau=\tfr15$; $\tau_1=\tfr3{25}$. Schedule: common control $w=(\tfr3{20},0)$ on $[0,\tau_1]$; in cell $\{1\}$, $u=-n_1$ from $\tau_1$ to rest; in cell $\{2,3\}$, the shared control $u=n_1$ on $[\tau_1,\tau]$ (the optimal shared brake of the pair, authority $\tfr35$), then $u=-n_j$ to rest, then $u=0$.

\paragraph{Branch 1.}
$n_1^\top w=\tfr3{20}>0$: the common control \emph{accelerates} branch~1 slightly, buying braking authority for the other two. At $\tau_1$: $\nu_1=1+\tfr3{20}\cdot\tfr3{25}=\tfr{509}{500}$, $\zeta_1=\tfr{34}{25}+\tfr3{25}+\tfrac12\cdot\tfr3{20}\cdot\tfr9{625}=\tfr{37027}{25000}$. Braking to rest from $\tau_1$ takes $\tfr{509}{500}<T-\tau_1=\tfr{27}{25}$ and adds $\tfrac12(\tfr{509}{500})^2=\tfr{259081}{500000}$: peak $\zeta_1^{\max}=\tfr{740540+259081}{500000}=\tfr{999621}{500000}=1.999242<2$. Velocity components: $(1,0)\to(\tfr{509}{500},0)\to(0,0)$; the largest magnitude $\tfr{509}{500}<\tfr65$.

\paragraph{Branches 2 and 3.}
$n_j^\top w=-\tfr9{100}$. At $\tau_1$: $\nu_j=1-\tfr9{100}\cdot\tfr3{25}=\tfr{2473}{2500}$, $\zeta_j=\tfr{34}{25}+\tfr3{25}-\tfrac12\cdot\tfr9{100}\cdot\tfr9{625}=\tfr{184919}{125000}$. Shared leg $u=n_1$ of length $\tfr2{25}$ with $n_j^\top n_1=-\tfr35$: $\nu_j\mapsto\tfr{2473}{2500}-\tfr6{125}=\tfr{2353}{2500}$, $\zeta_j\mapsto\tfr{184919}{125000}+\tfr{2473}{2500}\cdot\tfr2{25}-\tfrac12\cdot\tfr35\cdot\tfr4{625}=\tfr{194571}{125000}$. Braking to rest from $\tau$ takes $\tfr{2353}{2500}<1$ and adds $\tfrac12(\tfr{2353}{2500})^2=\tfr{5536609}{12500000}$: peak $\tfr{19457100+5536609}{12500000}=\tfr{24993709}{12500000}\approx1.999497<2$. Velocity components for branch~2: $(-\tfr35,\tfr45)\to(-\tfr35-\tfr9{500},\tfr45)\to(-\tfr35-\tfr9{500}+\tfr2{25},\tfr45)\to$ rest; all magnitudes at most $\tfr45$ (branch~3 is the mirror image). Noncritical position facets are below $\tfr{18}{25}$ by the bound of Section~\ref{S:pairs}.

\paragraph{Why these numbers.}
The common control is chosen on the bisector of branches~2 and~3 ($w_2=0$) so that both brake equally; the single free parameter $w_1$ trades branch~1's acceleration against the others' braking. At $w_1=\tfr3{20}$ the three peaks are $1.99924$ and $1.99950$: nearly balanced, which is what a schedule at the edge of viability should look like. The schedule is one feasible point, not an optimum, and the open gap $[\tfr3{25},\tfr{6-\sqrt{29}}5]$ is a statement about this schedule and the two-label dual jointly; closing it would require either a better schedule (pushing the lower end up) or a dual with more labels (pushing the upper end down). The finite program at $h=\tfr1{100}$ does not improve the upper end: its sweep (Table~\ref{tab:sweep}) first turns positive at $\tau_1=\tfr{13}{100}>\tau_1^{\mathrm{dual}}$, because of the moment error $Th/4=\tfr3{1000}$.

\begin{table}[h]
\centering\footnotesize
\caption{The archived sweep of Figure~4: finite program (10) at $h=\tfr1{100}$ for the two-stage structure, HiGHS optimum $\rho(\tau_1)$ and the eroded triple certificate $\Gamma(\tau_1)$. Every $\rho$ is a lower bound on $J$ for that structure; $\rho\ge\Gamma-Th/4$ throughout, with equality at $\tau_1=\tfr15$.}
\label{tab:sweep}
\begin{tabular}{@{}rrr@{\qquad}rrr@{}}
\toprule
$\tau_1$ & $\rho(\tau_1)$ & $\Gamma(\tau_1)$ & $\tau_1$ & $\rho(\tau_1)$ & $\Gamma(\tau_1)$\\
\midrule
$0.00$ & $-0.067500$ & $-0.105000$ & $0.11$ & $-0.011088$ & $-0.010537$\\
$0.01$ & $-0.067500$ & $-0.096037$ & $0.12$ & $-0.003450$ & $-0.002400$\\
$0.02$ & $-0.067500$ & $-0.087150$ & $0.13$ & $+0.004163$ & $+0.005662$\\
$0.03$ & $-0.067500$ & $-0.078338$ & $0.14$ & $+0.011775$ & $+0.013650$\\
$0.04$ & $-0.065095$ & $-0.069600$ & $0.15$ & $+0.019363$ & $+0.021562$\\
$0.05$ & $-0.057325$ & $-0.060937$ & $0.16$ & $+0.026925$ & $+0.029400$\\
$0.06$ & $-0.049575$ & $-0.052350$ & $0.17$ & $+0.034463$ & $+0.037163$\\
$0.07$ & $-0.041837$ & $-0.043838$ & $0.18$ & $+0.042000$ & $+0.044850$\\
$0.08$ & $-0.034125$ & $-0.035400$ & $0.19$ & $+0.049513$ & $+0.052463$\\
$0.09$ & $-0.026434$ & $-0.027037$ & $0.20$ & $+0.057000$ & $+0.060000$\\
$0.10$ & $-0.018750$ & $-0.018750$ & & &\\
\bottomrule
\end{tabular}
\end{table}

For $\tau_1\le\tfr3{100}$ the program's value is constant at $-0.0675$: there the binding sub-belief is the pair $\{2,3\}$ (branch~1 is revealed almost immediately and is harmless on its own), and $-0.0675$ is the finite value of that pair's program, consistent with the pair's exact slack $-\tfr{81}{1250}=-0.0648$ minus a moment error of at most $\tfr3{1000}$. This is the sense in which the optimised program is more informative than the fixed-weight dual: it re-optimises the weights and labels as the information structure changes, and its value switches from the triple's certificate to the pair's.
